\chapter{Security}
\label{ch:security}

The bot holds a server-admin token. Whoever controls the bot controls every account on the
server. The design therefore keeps that token in one place, limits who can make the bot do
anything, and makes bulk actions explicit.

\section{Who can make the bot act}

\begin{itemize}
  \item \textbf{Allowlist.} Only bot admins (instance config \cmd{admins:}), professors (DB
        allowlist, managed by admins) and course staff (DB, managed by course owners) can use
        the course commands. Other local users only get the public commands (\cmd{!help},
        \cmd{!password}, \cmd{!password code}) in a DM, which act on their own account and
        reveal nothing about courses or other people; everything else from them is ignored
        silently.
  \item \textbf{No implicit rights.} An email domain, the Synapse server-admin flag, the maubot
        UI login, and room power levels all grant nothing in the bot. Roles come only from
        the bot's configuration and database.
  \item \textbf{One permission check.} Every handler calls \cmd{authz.require(role, course)}.
        Admin satisfies everything; for a course, \emph{professor} means \enquote{owner of
        this course} and \emph{TA} means \enquote{owner or TA of this course}.
  \item \textbf{Only local senders in a real DM.} Commands are accepted only from users on this
        server (federation is off anyway) and only in a room with exactly the bot and the
        sender (\autoref{sec:dm-rule}), checked on every command. Links, email lists and
        tokens therefore never appear in a shared room.
  \item \textbf{Confirm for bulk and destructive actions.} Plans show what will happen; only
        the requester can confirm; the role is checked again at confirm time.
  \item \textbf{Resets for others are admin-only.} Whoever issues a reset link can open it. A
        professor issuing resets could take over students' accounts. The self-service paths
        (\cmd{!password}, reset codes) only ever work for the user's own account. Every reset
        is audited and the affected user is told in a DM (\autoref{sec:notices}).
\end{itemize}

\section{Signup and reset links}

\begin{itemize}
  \item Tokens are 32 random bytes (URL-safe). Only their SHA-256 hash is stored, so a
        database dump doesn't contain usable links.
  \item Single use: the token is burned atomically before the account is created, and
        restored only if account creation fails. Signup links expire after 14 days, reset
        links from an admin after 24 hours, links from \cmd{!password} after 1 hour.
  \item The username is fixed by the token. The page never lets the user choose it, so the
        mapping from university email to account can't be bypassed.
  \item The form carries a CSRF nonce: an HMAC of the token hash with a random key that
        exists only in memory.
  \item Rate limits per client IP and globally (\cmd{signup.rate_limit}), in addition to
        nginx's \cmd{limit_req}. The client IP is taken from \cmd{X-Forwarded-For} only if the
        connection comes from a configured proxy network (\cmd{signup.trusted_proxies}).
  \item The pages are minimal HTML without external assets, with a strict
        Content-Security-Policy, \cmd{no-store} caching, no referrer and no framing. Errors
        are generic and don't reveal whether a token existed.
  \item The minimum password length (\cmd{signup.min_password_length}, default 12) is enforced
        by the bot, because the admin API bypasses Synapse's password policy.
\end{itemize}

\section{Self-service password resets}

\begin{description}[style=nextline,leftmargin=1.5em,font=\normalfont\bfseries]
  \item[\cmd{!password}]
    The command has to come from a DM of the account itself, so the sender already holds a
    logged-in session. The link only works for that account. Because of that, the reset
    keeps all sessions (\cmd{logout_devices: false}); logging them out would only destroy
    the encryption keys of the person who asked. Someone who stole a session could also use
    it, but such a person can already read and write as the user; the user is notified and
    can take the account back with their reset code or an admin.
  \item[Reset codes]
    20 characters from Crockford's base32 alphabet (100 bits, no I, L, O, U), stored as a
    SHA-256 hash and compared in constant time. Guessing is pointless at that length; the
    limits are a second line of defence: the page's per-IP and global rate limit, and a
    per-account lock after \cmd{password.max_code_attempts} (default 5) wrong codes within
    24 hours, whatever the IP. The lock can be abused to block one person's code for a day;
    an admin reset still works then. Wrong code, unknown user and locked account all get the
    same answer. A used code is burned atomically before the password changes (restored if
    Synapse fails), all sessions are logged out (whoever uses the code may have lost them,
    or may be an attacker who must be kicked out), and a new code is shown.
  \item[What can't be fixed]
    Element's recovery key and the cross-signing private keys never leave the user's
    devices. No password reset can restore access to old encrypted messages. The server can
    only push users to save the key (well-known setting, I11), remind them in the welcome
    message, and explain \ui{Reset identity}.
\end{description}

\section{Professor bots}
\label{sec:why-userbots}

maubot plugins are ordinary Python in the maubot process, without a sandbox. A plugin can
import maubot's internals and read every client's access token. If professor plugins ran in
the main maubot, any uploaded plugin could take the admin bot's server-admin token.

\begin{itemize}
  \item \textbf{Separate maubot.} Professor bots run in \cmd{userbots}: its own container and
        database, no published port, no registration secret, and only the admin bot as its
        admin.
  \item \textbf{Startup guard.} The admin bot lists the target maubot's clients. If its own
        client appears there, the target is the main maubot, and bot commands are refused.
  \item \textbf{Package checks.} Size limits (packed and unpacked), \file{maubot.yaml} present
        and valid, the main module exists, no absolute or \cmd{..} paths, dependencies
        available in the pinned image. The SHA-256 of the package is recorded and shown to
        the approving admin.
  \item \textbf{Plugin ID ownership.} Plugin IDs are global in a maubot, and an upload with an
        existing ID replaces that plugin for every instance. The bot records the first
        uploader of each ID and refuses it for everyone else; the admin bot's own ID is
        reserved.
  \item \textbf{Approval, quota, kill switch.} Requests need an admin's approval (default),
        each professor has a quota (default 3), and \cmd{!admin bots disable-all} stops every
        professor bot.
  \item \textbf{Unprivileged accounts.} Bot accounts are normal users with \cmd{user_type: bot},
        named \cmd{@bot.<owner>.<name>}, and not exempt from rate limits unless configured.
  \item \textbf{Recovery keys.} Each bot's cross-signing recovery key is stored in the admin
        bot's database, which professor plugins can't reach, and is never shown to anyone.
        Whoever has it can sign devices as that bot account, so it is treated like the
        bot's access token: unprivileged, but kept out of logs and replies.
\end{itemize}

\textbf{Remaining risk:} professor bots share one process with \emph{each other}. A
malicious professor plugin could read another professor bot's token. Those tokens belong to
unprivileged accounts; approval and the kill switch limit the exposure.

\section{Known gaps}

\begin{itemize}
  \item \textbf{Open registration (I8).} As long as Synapse's token registration is enabled,
        someone with a registration token can pick any username, bypassing the email
        mapping. It should be switched off now that signup links are live.
  \item \textbf{Default crypto pickle key (I3)} on the main maubot (\autoref{ch:synapse}).
  \item \textbf{No offboarding.} Nothing tells the server when someone leaves the university;
        deactivation is manual.
  \item \textbf{Domain merge assumption.} \cmd{x@unibas.ch} and \cmd{x@stud.unibas.ch} are
        treated as one person. This is correct only if the university never gives the same
        name to two people on the two domains (to be confirmed with ITS).
  \item \textbf{Reset codes in chat history.} The reply to \cmd{!password code} stays in the
        user's encrypted DM with the bot. Anyone who can read that DM can already act as the
        user. Users are told to keep the code outside Matrix too.
  \item \textbf{Not yet verified on the test node:} encrypted uploads, the real nginx route,
        and that held notices are readable on the user's first device.
\end{itemize}
